\section{Appendix: sentai.io}\label{appendix-sentai.io}

The \texttt{sentai.io} namespace groups the simplest board-local control
primitives. In the current runtime it is intentionally small and is used
mainly for visual status indication, quick feedback during debugging,
and minimal interaction from scripts that do not require a richer GPIO
abstraction.

\subsection{Functions}\label{functions}

\begin{itemize}
\tightlist
\item
  \texttt{sentai.io.led\_on()}: Turns the user LED on.
\item
  \texttt{sentai.io.led\_off()}: Turns the user LED off.
\end{itemize}

\subsection{Example}\label{example}

\begin{verbatim}
import sentai

for _ in range(3):
    sentai.io.led_on()
    sentai.rtos.sleep_ms(200)
    sentai.io.led_off()
    sentai.rtos.sleep_ms(200)
\end{verbatim}

\section{Appendix: sentai.rtos}\label{appendix-sentai.rtos}

The \texttt{sentai.rtos} namespace exposes runtime observability and
timing services from the FreeRTOS layer. It is useful both for
interactive diagnosis and for scripts that need coarse scheduling,
uptime measurement, or visibility into task and memory behavior during
long-running onboard experiments.

\subsection{Functions}\label{functions-1}

\begin{itemize}
\tightlist
\item
  \texttt{sentai.rtos.sleep\_ms(ms)}: Suspends the current script for a
  specified number of milliseconds.
\item
  \texttt{sentai.rtos.ticks\_ms()}: Returns a monotonic millisecond
  counter.
\item
  \texttt{sentai.rtos.tasks()}: Returns the current task list with
  state, priority, and stack high-water mark.
\item
  \texttt{sentai.rtos.heap\_info()}: Returns memory statistics for the
  RTOS heap and MicroPython GC heap.
\item
  \texttt{sentai.rtos.cpu\_usage()}: Returns per-task CPU-usage
  estimates.
\item
  \texttt{sentai.rtos.uptime()}: Returns system uptime in seconds.
\end{itemize}

\subsection{Example}\label{example-1}

\begin{verbatim}
import sentai

print('uptime:', sentai.rtos.uptime())
print('heap:', sentai.rtos.heap_info())
for name, state, prio, hwm in sentai.rtos.tasks():
    print(name, state, prio, hwm)
\end{verbatim}

\section{Appendix: sentai.tpu}\label{appendix-sentai.tpu}

The \texttt{sentai.tpu} namespace is the main interface to
EdgeTPU-backed inference. It covers model loading, input preparation,
execution, tensor inspection, quantization helpers, and
detector-oriented post-processing. In practice it is the default path
for compact vision models compiled for the accelerator.

\subsection{Functions}\label{functions-2}

\begin{itemize}
\tightlist
\item
  \texttt{sentai.tpu.load(path)}: Loads a \texttt{.tflite} model from
  flash and creates the EdgeTPU interpreter.
\item
  \texttt{sentai.tpu.load\_image(path)}: Loads a JPEG or raw RGB image
  into the input tensor.
\item
  \texttt{sentai.tpu.invoke()}: Runs inference and returns elapsed time
  in milliseconds.
\item
  \texttt{sentai.tpu.ready()}: Reports whether the model and interpreter
  are ready.
\item
  \texttt{sentai.tpu.num\_outputs()}: Returns the number of output
  tensors.
\item
  \texttt{sentai.tpu.output\_size(idx)}: Returns the byte size of an
  output tensor.
\item
  \texttt{sentai.tpu.output(idx)}: Returns raw output bytes.
\item
  \texttt{sentai.tpu.output\_dims(idx)}: Returns the output tensor
  shape.
\item
  \texttt{sentai.tpu.output\_type(idx)}: Returns the tensor type.
\item
  \texttt{sentai.tpu.row(idx,\ r)}: Returns a selected tensor row.
\item
  \texttt{sentai.tpu.value(idx,\ flat\_i)}: Returns a single flattened
  value.
\item
  \texttt{sentai.tpu.save\_output(path)}: Saves all outputs to CSV.
\item
  \texttt{sentai.tpu.input\_quant()}: Returns input quantization
  parameters.
\item
  \texttt{sentai.tpu.output\_quant(idx)}: Returns output quantization
  parameters.
\item
  \texttt{sentai.tpu.output\_floats(idx)}: Returns dequantized outputs
  as Python floats.
\item
  \texttt{sentai.tpu.input\_type()}: Returns the input tensor type.
\item
  \texttt{sentai.tpu.detect(conf,\ iou,\ max)}: Runs YOLO-style
  non-maximum suppression on the last inference output.
\item
  \texttt{sentai.tpu.draw(path,\ dets,\ quality)}: Draws detections on
  the last captured frame and saves a JPEG.
\end{itemize}

\subsection{Example}\label{example-2}

\begin{verbatim}
import sentai

sentai.tpu.load('/models/yolov8n.tflite')
sentai.camera.init()
sentai.camera.to_tensor()
ms = sentai.tpu.invoke()
dets = sentai.tpu.detect(0.25, 0.45, 50)
print('ms:', ms, 'detections:', len(dets))
sentai.tpu.draw('/det.jpg', dets)
\end{verbatim}

\section{Appendix: sentai.tfl}\label{appendix-sentai.tfl}

The \texttt{sentai.tfl} namespace provides CPU-side TensorFlow Lite
Micro inference on the Cortex-M7. It complements \texttt{sentai.tpu} for
models that are not compiled for the EdgeTPU, for smaller auxiliary
models, or for experiments that require direct MCU execution with
explicit tensor-arena sizing.

\subsection{Functions}\label{functions-3}

\begin{itemize}
\tightlist
\item
  \texttt{sentai.tfl.load(path,\ arena\_kb)}: Loads a \texttt{.tflite}
  model and allocates the tensor arena.
\item
  \texttt{sentai.tfl.unload()}: Frees the interpreter and tensor arena.
\item
  \texttt{sentai.tfl.invoke()}: Runs inference and returns elapsed time
  in milliseconds.
\item
  \texttt{sentai.tfl.ready()}: Reports whether the interpreter is ready.
\item
  \texttt{sentai.tfl.info()}: Prints model and arena information.
\item
  \texttt{sentai.tfl.set\_input(data)}: Writes bytes into the input
  tensor.
\item
  \texttt{sentai.tfl.load\_image(path)}: Loads a JPEG or raw RGB image
  into the input tensor.
\item
  \texttt{sentai.tfl.input\_size()}: Returns the input tensor size in
  bytes.
\item
  \texttt{sentai.tfl.input\_dims()}: Returns the input tensor shape.
\item
  \texttt{sentai.tfl.input\_type()}: Returns the input tensor type.
\item
  \texttt{sentai.tfl.input\_quant()}: Returns input quantization
  parameters.
\item
  \texttt{sentai.tfl.num\_outputs()}: Returns the number of output
  tensors.
\item
  \texttt{sentai.tfl.output\_size(idx)}: Returns output tensor size.
\item
  \texttt{sentai.tfl.output(idx)}: Returns raw output bytes.
\item
  \texttt{sentai.tfl.output\_dims(idx)}: Returns the output tensor
  shape.
\item
  \texttt{sentai.tfl.output\_type(idx)}: Returns the output tensor type.
\item
  \texttt{sentai.tfl.output\_quant(idx)}: Returns output quantization
  parameters.
\item
  \texttt{sentai.tfl.output\_floats(idx)}: Returns dequantized output
  values.
\item
  \texttt{sentai.tfl.row(idx,\ r)}: Returns one tensor row.
\item
  \texttt{sentai.tfl.value(idx,\ flat\_i)}: Returns one output value.
\item
  \texttt{sentai.tfl.save\_output(path)}: Saves all outputs to CSV.
\end{itemize}

\subsection{Example}\label{example-3}

\begin{verbatim}
import sentai

sentai.tfl.load('/models/keyword_detect.tflite', 256)
sentai.tfl.set_input(feature_bytes)
ms = sentai.tfl.invoke()
probs = sentai.tfl.output_floats(0)
print('inference:', ms, 'ms', probs)
\end{verbatim}

\section{Appendix: sentai.fs}\label{appendix-sentai.fs}

The \texttt{sentai.fs} namespace provides script-level access to the
LittleFS user partition. It is used to move models, scripts, logs,
images, and outputs between the runtime and persistent flash, and it is
a central part of the interactive workflow because both model deployment
and experiment logging depend on it.

\subsection{Functions}\label{functions-4}

\begin{itemize}
\tightlist
\item
  \texttt{sentai.fs.read(path)}: Reads an entire file as bytes.
\item
  \texttt{sentai.fs.read\_str(path)}: Reads an entire file as text.
\item
  \texttt{sentai.fs.read\_base64(path)}: Reads and prints a file as
  base64 text.
\item
  \texttt{sentai.fs.write(path,\ data)}: Writes bytes or text to flash.
\item
  \texttt{sentai.fs.size(path)}: Returns the file size.
\item
  \texttt{sentai.fs.exists(path)}: Checks whether a file or directory
  exists.
\item
  \texttt{sentai.fs.remove(path)}: Deletes a file or empty directory.
\item
  \texttt{sentai.fs.mkdir(path)}: Creates directories recursively.
\item
  \texttt{sentai.fs.ls(path)}: Lists directory contents.
\item
  \texttt{sentai.fs.format()}: Reformats the user partition and erases
  user data.
\end{itemize}

\subsection{Example}\label{example-4}

\begin{verbatim}
import sentai

sentai.fs.mkdir('/logs')
sentai.fs.write('/logs/run.txt', 'experiment started\n')
print(sentai.fs.read_str('/logs/run.txt'))
print(sentai.fs.ls('/logs'))
\end{verbatim}

\section{Appendix: sentai.camera}\label{appendix-sentai.camera}

The \texttt{sentai.camera} namespace controls image capture and basic
camera-side preprocessing. It provides the main bridge between the
sensing layer and the inference pipeline, including frame capture, JPEG
export, tensor feeding, camera selection, and frame-sequence inspection
for dual-camera operation.

\subsection{Functions}\label{functions-5}

\begin{itemize}
\tightlist
\item
  \texttt{sentai.camera.init(streaming)}: Starts the camera in
  continuous or trigger mode.
\item
  \texttt{sentai.camera.stop()}: Stops camera capture.
\item
  \texttt{sentai.camera.jpeg(quality)}: Captures a frame and returns
  JPEG bytes.
\item
  \texttt{sentai.camera.save\_jpeg(path,\ q)}: Captures and saves a JPEG
  file.
\item
  \texttt{sentai.camera.to\_tensor()}: Captures a frame and writes it
  into the TPU input tensor.
\item
  \texttt{sentai.camera.resolution()}: Returns the current output
  resolution.
\item
  \texttt{sentai.camera.set\_resolution(w,\ h)}: Changes the capture
  resolution.
\item
  \texttt{sentai.camera.native\_res()}: Returns the sensor native
  resolution.
\item
  \texttt{sentai.camera.select(id)}: Selects the active camera.
\item
  \texttt{sentai.camera.frame\_count()}: Returns the monotonic hardware
  frame counter.
\end{itemize}

\subsection{Example}\label{example-5}

\begin{verbatim}
import sentai

sentai.camera.init()
sentai.camera.set_resolution(320, 320)
sentai.camera.select(0)
sentai.camera.save_jpeg('/front.jpg', 85)
sentai.camera.select(1)
sentai.camera.save_jpeg('/back.jpg', 85)
sentai.camera.stop()
\end{verbatim}

\section{Appendix: sentai.imu}\label{appendix-sentai.imu}

The \texttt{sentai.imu} namespace exposes the onboard LIS2DU12
accelerometer. It provides direct access to acceleration and derived
tilt estimates and is commonly used for motion monitoring, camera-pose
estimation, and simple event detection in scripts that combine
perception with inertial context.

\subsection{Functions}\label{functions-6}

\begin{itemize}
\tightlist
\item
  \texttt{sentai.imu.init()}: Initializes the accelerometer.
\item
  \texttt{sentai.imu.read()}: Returns acceleration and temperature
  measurements.
\item
  \texttt{sentai.imu.degrees()}: Returns pitch and roll in degrees.
\item
  \texttt{sentai.imu.radians()}: Returns pitch and roll in radians.
\end{itemize}

\subsection{Example}\label{example-6}

\begin{verbatim}
import sentai

sentai.imu.init()
for _ in range(10):
    d = sentai.imu.degrees()
    if d:
        print('pitch:', d['pitch'], 'roll:', d['roll'])
    sentai.rtos.sleep_ms(100)
\end{verbatim}

\section{Appendix: sentai.mic}\label{appendix-sentai.mic}

The \texttt{sentai.mic} namespace manages the onboard PDM microphone
through a ring-buffer recording model. It is designed for low-overhead
acquisition, sound-level monitoring, and asynchronous MP3 export,
allowing audio-triggered experiments without forcing the user to manage
raw DMA buffers directly.

\subsection{Functions}\label{functions-7}

\begin{itemize}
\tightlist
\item
  \texttt{sentai.mic.start(seconds)}: Starts ring-buffer recording for
  the specified window length.
\item
  \texttt{sentai.mic.stop()}: Stops the microphone and returns the
  number of buffered samples.
\item
  \texttt{sentai.mic.recording()}: Reports whether the microphone is
  active.
\item
  \texttt{sentai.mic.samples()}: Returns the number of samples currently
  available.
\item
  \texttt{sentai.mic.level()}: Returns the current RMS level in
  centi-decibels.
\item
  \texttt{sentai.mic.save\_mp3()}: Saves buffered audio as an MP3 file
  and returns the filename.
\end{itemize}

\subsection{Example}\label{example-7}

\begin{verbatim}
import sentai

sentai.mic.start(5)
sentai.rtos.sleep_ms(5000)
name = sentai.mic.save_mp3()
print('saved:', name)
sentai.mic.stop()
\end{verbatim}

\section{Appendix: sentai.usb}\label{appendix-sentai.usb}

The \texttt{sentai.usb} namespace provides two distinct services:
mass-storage export of the user partition and raw USB CDC serial I/O
when the REPL has been moved away from USB. It is therefore both a
deployment path for files and a host-device communication channel, with
explicit coordination required to avoid conflicts with filesystem
access.

\subsection{Functions}\label{functions-8}

\begin{itemize}
\tightlist
\item
  \texttt{sentai.usb.drive(on)}: Enables or disables USB mass-storage
  export.
\item
  \texttt{sentai.usb.open()}: Opens USB CDC ACM for raw serial I/O.
\item
  \texttt{sentai.usb.close()}: Closes the raw USB CDC channel.
\item
  \texttt{sentai.usb.write(data)}: Writes bytes or text to the USB host.
\item
  \texttt{sentai.usb.read(max,\ timeout\_ms)}: Reads bytes from the USB
  host.
\item
  \texttt{sentai.usb.available()}: Returns the number of pending
  received bytes.
\end{itemize}

\subsection{Example}\label{example-8}

\begin{verbatim}
import sentai

sentai.usb.drive(1)
# host copies files here
sentai.usb.drive(0)
print(sentai.fs.ls('/'))
\end{verbatim}

\section{Appendix: sentai.uart}\label{appendix-sentai.uart}

The \texttt{sentai.uart} namespace exposes raw UART serial I/O on the
board-side serial port when that port is not reserved for the REPL or
another protocol bridge. It is intended for direct communication with
external serial devices such as radios, sensors, GPS modules, or custom
peripherals.

\subsection{Functions}\label{functions-9}

\begin{itemize}
\tightlist
\item
  \texttt{sentai.uart.open(baud)}: Opens the UART port at the requested
  baud rate.
\item
  \texttt{sentai.uart.close()}: Closes the UART port and restores the
  default configuration.
\item
  \texttt{sentai.uart.write(data)}: Sends bytes or text.
\item
  \texttt{sentai.uart.read(max,\ timeout\_ms)}: Reads bytes with
  polling, timeout, or blocking behavior.
\item
  \texttt{sentai.uart.available()}: Returns the number of waiting bytes.
\end{itemize}

\subsection{Example}\label{example-9}

\begin{verbatim}
import sentai

sentai.uart.open(38400)
sentai.uart.write(b'AT\r\n')
resp = sentai.uart.read(256, 1000)
print(resp)
sentai.uart.close()
\end{verbatim}

\section{Appendix: sentai.console}\label{appendix-sentai.console}

The \texttt{sentai.console} interface controls where the interactive
MicroPython REPL is exposed. Although small, it is operationally
important because several communication namespaces reuse the non-REPL
interface, so moving the console between USB and UART determines which
transport remains available for scripting.

\subsection{Functions}\label{functions-10}

\begin{itemize}
\tightlist
\item
  \texttt{sentai.console()}: Returns the current REPL target, either
  \texttt{usb} or \texttt{uart}.
\item
  \texttt{sentai.console(\textquotesingle{}usb\textquotesingle{})}:
  Moves the REPL to USB CDC ACM.
\item
  \texttt{sentai.console(\textquotesingle{}uart\textquotesingle{})}:
  Moves the REPL to the UART console.
\end{itemize}

\subsection{Example}\label{example-10}

\begin{verbatim}
import sentai

print(sentai.console())
sentai.console('uart')
sentai.usb.open()
sentai.usb.write(b'host channel ready\n')
sentai.usb.close()
\end{verbatim}

\section{Appendix: sentai.mesh}\label{appendix-sentai.mesh}

The \texttt{sentai.mesh} namespace bridges the runtime to
Meshtastic-compatible radios over UART. It supports both simple text
exchange and structured vision-message transport, making it suitable for
distributed sensing, low-bandwidth telemetry, and multi-node experiments
in which detections or updates must be serialized and disseminated over
a mesh network.

\subsection{Functions}\label{functions-11}

\begin{itemize}
\tightlist
\item
  \texttt{sentai.mesh.init(baud)}: Opens the UART link and starts the
  mesh receive task.
\item
  \texttt{sentai.mesh.stop()}: Stops the receive task and closes the
  link.
\item
  \texttt{sentai.mesh.node()}: Returns the local node number.
\item
  \texttt{sentai.mesh.config()}: Requests configuration and NodeDB
  information from the radio.
\item
  \texttt{sentai.mesh.send(text,\ dest,\ ch,\ ack)}: Sends a text
  message.
\item
  \texttt{sentai.mesh.available()}: Returns the number of queued text
  messages.
\item
  \texttt{sentai.mesh.receive(timeout)}: Receives the next text message.
\item
  \texttt{sentai.mesh.send\_detection(...)}: Sends a structured
  new-detection message.
\item
  \texttt{sentai.mesh.send\_update(...)}: Sends a structured
  track-update message.
\item
  \texttt{sentai.mesh.vision\_available()}: Returns the number of queued
  vision messages.
\item
  \texttt{sentai.mesh.receive\_vision(timeout)}: Receives the next
  structured vision message.
\item
  \texttt{sentai.mesh.set\_pose(pitch\_deg,\ roll\_deg,\ altitude\_cm,\ heading\_deg)}:
  Attaches camera pose metadata to outgoing vision messages.
\end{itemize}

\subsection{Example}\label{example-11}

\begin{verbatim}
import sentai

sentai.mesh.init()
sentai.mesh.send('hello world')
msg = sentai.mesh.receive(5000)
if msg:
    print(msg['from'], msg['text'])
\end{verbatim}

\section{Appendix: sentai.link}\label{appendix-sentai.link}

The \texttt{sentai.link} namespace exposes a MAVLink v2 bridge over
UART. It supports message reception, heartbeat and status transmission,
command dispatch, structured vision reporting, and obstacle-map
generation for PX4 collision prevention. It is the primary interface for
integrating the runtime with a conventional autopilot stack.

\subsection{Functions}\label{functions-12}

\begin{itemize}
\tightlist
\item
  \texttt{sentai.link.init(baud,\ sysid,\ compid)}: Opens the MAVLink
  link and starts the receive task.
\item
  \texttt{sentai.link.stop()}: Stops the link and closes the UART
  channel.
\item
  \texttt{sentai.link.available()}: Returns the number of queued MAVLink
  messages.
\item
  \texttt{sentai.link.receive(timeout)}: Returns the next decoded
  MAVLink message.
\item
  \texttt{sentai.link.heartbeat(type)}: Sends a heartbeat.
\item
  \texttt{sentai.link.send(text,\ sev)}: Sends a \texttt{STATUSTEXT}
  message.
\item
  \texttt{sentai.link.command(tsys,\ tcomp,\ cmd,\ conf,\ p1..p7)}:
  Sends a \texttt{COMMAND\_LONG} message.
\item
  \texttt{sentai.link.send\_detection(...)}: Sends a structured vision
  detection.
\item
  \texttt{sentai.link.send\_update(...)}: Sends a structured vision
  update.
\item
  \texttt{sentai.link.send\_delete(...)}: Sends a structured vision
  deletion event.
\item
  \texttt{sentai.link.obstacles\_from\_tracker(...)}: Builds and sends a
  72-bin obstacle map from active tracks.
\item
  \texttt{sentai.link.obstacle\_distance(distances72,\ ...)}: Sends a
  raw \texttt{OBSTACLE\_DISTANCE} message.
\item
  \texttt{sentai.link.obstacles\_from\_points(points,\ ...)}: Builds an
  obstacle map from arbitrary XY points.
\end{itemize}

\subsection{Example}\label{example-12}

\begin{verbatim}
import sentai

sentai.pipeline.init()
sentai.pipeline.set_pose(120, -1, 0.0, 0.0)
sentai.link.init(57600)
while True:
    sentai.link.obstacles_from_tracker(800, 20)
    sentai.link.heartbeat()
    sentai.rtos.sleep_ms(100)
\end{verbatim}

\section{Appendix: sentai.crazy}\label{appendix-sentai.crazy}

The \texttt{sentai.crazy} namespace connects the runtime to a Crazyflie
platform over CPX/CRTP. It exposes both higher-level flight procedures
and lower-level control interfaces, allowing scripted indoor flight
experiments without moving mission logic outside the main MicroPython
environment.

\subsection{Functions}\label{functions-13}

\begin{itemize}
\tightlist
\item
  \texttt{sentai.crazy.init(baud)}: Starts the Crazyflie communication
  tasks.
\item
  \texttt{sentai.crazy.stop()}: Stops the bridge and closes the link.
\item
  \texttt{sentai.crazy.debug(level)}: Changes debug verbosity.
\item
  \texttt{sentai.crazy.arm()}: Arms the drone.
\item
  \texttt{sentai.crazy.disarm()}: Disarms the drone.
\item
  \texttt{sentai.crazy.ping(timeout)}: Measures round-trip time to the
  drone.
\item
  \texttt{sentai.crazy.fly(height,\ hold\_ms,\ takeoff\_ms,\ land\_ms)}:
  Executes a blocking high-level flight cycle.
\item
  \texttt{sentai.crazy.attitude(roll,\ pitch,\ yaw\_rate,\ thrust)}:
  Sends non-blocking attitude setpoints.
\item
  \texttt{sentai.crazy.fly\_stop()}: Stops attitude-controlled flight
  and disarms.
\item
  \texttt{sentai.crazy.takeoff(h,\ dur,\ yaw,\ use\_yaw,\ group)}: Sends
  a high-level takeoff command.
\item
  \texttt{sentai.crazy.land(h,\ dur,\ yaw,\ use\_yaw,\ group)}: Sends a
  high-level landing command.
\item
  \texttt{sentai.crazy.stop\_motors(group)}: Performs an emergency motor
  stop.
\item
  \texttt{sentai.crazy.go\_to(x,\ y,\ z,\ yaw,\ dur,\ rel,\ lin,\ group)}:
  Sends a waypoint-style movement command.
\item
  \texttt{sentai.crazy.hover(vx,\ vy,\ yr,\ z)}: Sends hover-mode
  velocity commands.
\item
  \texttt{sentai.crazy.test\_fly(power,\ dur\_ms)}: Runs a raw motor
  test.
\item
  \texttt{sentai.crazy.send\_crtp(port,\ ch,\ data)}: Sends an arbitrary
  CRTP packet.
\end{itemize}

\subsection{Example}\label{example-13}

\begin{verbatim}
import sentai

sentai.crazy.init()
sentai.crazy.arm()
sentai.crazy.takeoff(0.5, 2.0)
sentai.rtos.sleep_ms(3000)
sentai.crazy.land(0.0, 2.0)
sentai.crazy.stop()
\end{verbatim}

\section{Appendix: sentai.pipeline}\label{appendix-sentai.pipeline}

The \texttt{sentai.pipeline} namespace exposes the background detection
pipeline and the SentAI-SORT tracker. It is the main high-level
perception service in the runtime and packages detector scheduling,
tracking, event reporting, camera geometry, and ground-plane projection
into a single script-facing interface.

\subsection{Functions}\label{functions-14}

\begin{itemize}
\tightlist
\item
  \texttt{sentai.pipeline.start(conf,\ iou,\ max,\ track)}: Starts
  continuous detection and optionally tracking.
\item
  \texttt{sentai.pipeline.stop()}: Stops the pipeline.
\item
  \texttt{sentai.pipeline.running()}: Reports whether the pipeline is
  active.
\item
  \texttt{sentai.pipeline.get(timeout\_ms)}: Returns the next detection
  frame.
\item
  \texttt{sentai.pipeline.stats()}: Returns processed, dropped, and FPS
  statistics.
\item
  \texttt{sentai.pipeline.tracks()}: Returns the current active-track
  snapshot.
\item
  \texttt{sentai.pipeline.event(timeout\_ms)}: Returns the next tracker
  event.
\item
  \texttt{sentai.pipeline.track\_config(...)}: Gets or sets tracker
  parameters.
\item
  \texttt{sentai.pipeline.set\_pose(altitude\_cm,\ heading\_deg,\ lat,\ lon)}:
  Sets camera pose for ground projection.
\item
  \texttt{sentai.pipeline.camera\_config(cam\_id,\ fov\_h,\ fov\_v,\ mount\_pitch,\ mount\_roll,\ mount\_yaw)}:
  Gets or sets per-camera geometry.
\end{itemize}

\subsection{Example}\label{example-14}

\begin{verbatim}
import sentai

sentai.camera.init()
sentai.pipeline.camera_config(0, 70.8, 43.4, 0, 0, 0)
sentai.pipeline.set_pose(1000, 0, 44.4268, 26.1025)
sentai.pipeline.start(0.3, 0.45, 50, True)
while True:
    evt = sentai.pipeline.event(2000)
    if evt:
        print(evt)
\end{verbatim}

\section{Appendix: sentai.aifes}\label{appendix-sentai.aifes}

The \texttt{sentai.aifes} namespace integrates the AIfES library for
on-device neural-network training and inference. It is intended for
small multilayer perceptrons and compact transfer-learning workflows in
which features extracted by the TPU or camera are fed into a trainable
MCU-side head.

\subsection{Functions}\label{functions-15}

\begin{itemize}
\tightlist
\item
  \texttt{sentai.aifes.load(path)}: Loads a model architecture from
  YAML.
\item
  \texttt{sentai.aifes.load\_weights(path)}: Loads pre-trained weights.
\item
  \texttt{sentai.aifes.save\_weights(path)}: Saves trained weights.
\item
  \texttt{sentai.aifes.unload()}: Frees the loaded model.
\item
  \texttt{sentai.aifes.ready()}: Reports whether the model is ready.
\item
  \texttt{sentai.aifes.train(x\_data,\ y\_data,\ ...)}: Trains the model
  on supplied samples.
\item
  \texttt{sentai.aifes.set\_input(data)}: Sets the input vector from
  floats.
\item
  \texttt{sentai.aifes.from\_tpu(idx)}: Uses a TPU output tensor as
  AIfES input.
\item
  \texttt{sentai.aifes.from\_camera(w,\ h,\ gray)}: Uses a camera frame
  as normalized input.
\item
  \texttt{sentai.aifes.from\_mic(samples)}: Uses microphone samples as
  input.
\item
  \texttt{sentai.aifes.invoke()}: Runs inference.
\item
  \texttt{sentai.aifes.output()}: Returns the latest output vector.
\item
  \texttt{sentai.aifes.mic\_init()}: Initializes microphone capture for
  audio models.
\item
  \texttt{sentai.aifes.mic\_stop()}: Stops microphone capture.
\item
  \texttt{sentai.aifes.predict(input)}: Legacy single-call inference
  helper.
\item
  \texttt{sentai.aifes.info()}: Returns model metadata.
\end{itemize}

\subsection{Example}\label{example-15}

\begin{verbatim}
import sentai

sentai.tpu.load('/models/mobilenet_features.tflite')
sentai.aifes.load('/models/classifier_head.yaml')
sentai.camera.init()

sentai.camera.to_tensor()
sentai.tpu.invoke()
sentai.aifes.from_tpu(0)
sentai.aifes.invoke()
print(sentai.aifes.output())
\end{verbatim}

\section{Appendix: sentai.kmeans}\label{appendix-sentai.kmeans}

The \texttt{sentai.kmeans} namespace implements compact k-means
clustering for embedding vectors. It is mainly useful for few-shot or
unsupervised classification workflows in which a TPU feature extractor
produces dense embeddings and the runtime clusters them directly on
device.

\subsection{Functions}\label{functions-16}

\begin{itemize}
\tightlist
\item
  \texttt{sentai.kmeans.init(k,\ dim)}: Initializes the clustering
  state.
\item
  \texttt{sentai.kmeans.add(cluster\_idx,\ vector)}: Adds a labeled
  vector to a cluster.
\item
  \texttt{sentai.kmeans.add\_from\_tpu(cluster\_idx,\ tpu\_idx)}: Adds a
  TPU output directly.
\item
  \texttt{sentai.kmeans.compute()}: Computes centroids from labeled
  samples.
\item
  \texttt{sentai.kmeans.fit(max\_iter)}: Runs unsupervised Lloyd
  iterations.
\item
  \texttt{sentai.kmeans.predict(vector)}: Predicts the closest centroid.
\item
  \texttt{sentai.kmeans.from\_tpu(tpu\_idx)}: Predicts from a TPU output
  tensor.
\item
  \texttt{sentai.kmeans.distances(vector)}: Returns distances to all
  centroids.
\item
  \texttt{sentai.kmeans.save(path)}: Saves centroids to flash.
\item
  \texttt{sentai.kmeans.load(path)}: Loads centroids from flash.
\item
  \texttt{sentai.kmeans.centroid(idx)}: Returns one centroid vector.
\item
  \texttt{sentai.kmeans.info()}: Returns configuration and sample-count
  metadata.
\item
  \texttt{sentai.kmeans.clear()}: Clears training vectors while keeping
  centroids.
\end{itemize}

\subsection{Example}\label{example-16}

\begin{verbatim}
import sentai

sentai.tpu.load('/models/mobilenet_features.tflite')
sentai.camera.init()
sentai.kmeans.init(3, 1280)

sentai.camera.to_tensor()
sentai.tpu.invoke()
sentai.kmeans.add_from_tpu(0, 0)
sentai.kmeans.compute()
print(sentai.kmeans.info())
\end{verbatim}

\section{Appendix: sentai.pca}\label{appendix-sentai.pca}

The \texttt{sentai.pca} namespace provides principal-component analysis
for high-dimensional embeddings. Its role is to reduce vector dimension
before downstream use in classifiers, clustering methods, or lightweight
anomaly detectors, especially when raw feature vectors are too large for
repeated on-device processing.

\subsection{Functions}\label{functions-17}

\begin{itemize}
\tightlist
\item
  \texttt{sentai.pca.init(in\_dim,\ out\_dim)}: Initializes the PCA
  model.
\item
  \texttt{sentai.pca.add(vector)}: Adds a training vector.
\item
  \texttt{sentai.pca.from\_tpu(idx)}: Adds a TPU output tensor as
  training data.
\item
  \texttt{sentai.pca.fit()}: Computes the principal components.
\item
  \texttt{sentai.pca.transform(vector)}: Projects a vector into the
  reduced space.
\item
  \texttt{sentai.pca.transform\_tpu(idx)}: Projects a TPU output tensor
  directly.
\item
  \texttt{sentai.pca.inverse(reduced)}: Reconstructs an approximate
  original vector.
\item
  \texttt{sentai.pca.explained\_variance()}: Returns component
  variances.
\item
  \texttt{sentai.pca.save(path)}: Saves the fitted PCA model.
\item
  \texttt{sentai.pca.load(path)}: Loads a PCA model.
\item
  \texttt{sentai.pca.info()}: Returns model metadata.
\item
  \texttt{sentai.pca.clear()}: Clears stored training vectors.
\end{itemize}

\subsection{Example}\label{example-17}

\begin{verbatim}
import sentai

sentai.tpu.load('/models/mobilenet_features.tflite')
sentai.camera.init()
sentai.pca.init(1280, 32)
sentai.camera.to_tensor()
sentai.tpu.invoke()
sentai.pca.from_tpu(0)
sentai.pca.fit()
print(sentai.pca.transform_tpu(0)[:5])
\end{verbatim}

\section{Appendix: sentai.anomaly}\label{appendix-sentai.anomaly}

The \texttt{sentai.anomaly} namespace implements online anomaly
detection based on incremental covariance estimation and Mahalanobis
distance, together with a simple CUSUM change detector. It is suitable
for learning a baseline normal distribution and then scoring embeddings
or scalar signals during deployment.

\subsection{Functions}\label{functions-18}

\begin{itemize}
\tightlist
\item
  \texttt{sentai.anomaly.init(dim)}: Initializes the anomaly model.
\item
  \texttt{sentai.anomaly.observe(vector)}: Updates the normal model with
  one observation.
\item
  \texttt{sentai.anomaly.observe\_tpu(idx)}: Updates the model from a
  TPU output tensor.
\item
  \texttt{sentai.anomaly.score(vector)}: Returns the anomaly score for a
  vector.
\item
  \texttt{sentai.anomaly.score\_tpu(idx)}: Scores a TPU output tensor.
\item
  \texttt{sentai.anomaly.threshold(val)}: Gets or sets the anomaly
  threshold.
\item
  \texttt{sentai.anomaly.is\_anomaly(vector)}: Tests whether a vector
  exceeds the threshold.
\item
  \texttt{sentai.anomaly.is\_anomaly\_tpu(idx)}: Tests a TPU output
  tensor.
\item
  \texttt{sentai.anomaly.cusum\_init(threshold,\ drift)}: Initializes
  the CUSUM detector.
\item
  \texttt{sentai.anomaly.cusum\_observe(value)}: Feeds one scalar sample
  into CUSUM.
\item
  \texttt{sentai.anomaly.cusum\_score()}: Returns current positive and
  negative cumulative sums.
\item
  \texttt{sentai.anomaly.cusum\_reset()}: Clears CUSUM accumulators.
\item
  \texttt{sentai.anomaly.save(path)}: Saves the learned model.
\item
  \texttt{sentai.anomaly.load(path)}: Loads a saved model.
\item
  \texttt{sentai.anomaly.stats()}: Returns model statistics.
\item
  \texttt{sentai.anomaly.info()}: Alias for \texttt{stats()}.
\item
  \texttt{sentai.anomaly.clear()}: Clears all statistics.
\end{itemize}

\subsection{Example}\label{example-18}

\begin{verbatim}
import sentai

sentai.tpu.load('/models/mobilenet_features.tflite')
sentai.camera.init()
sentai.anomaly.init(1280)

for _ in range(30):
    sentai.camera.to_tensor()
    sentai.tpu.invoke()
    sentai.anomaly.observe_tpu(0)

sentai.camera.to_tensor()
sentai.tpu.invoke()
print(sentai.anomaly.score_tpu(0))
\end{verbatim}

\section{Appendix: sentai.dtw}\label{appendix-sentai.dtw}

The \texttt{sentai.dtw} namespace exposes dynamic time warping for
template-based sequence matching. It is intended for short gesture or
audio patterns where full trainable sequence models are unnecessary and
a small number of stored templates is sufficient for online recognition.

\subsection{Functions}\label{functions-19}

\begin{itemize}
\tightlist
\item
  \texttt{sentai.dtw.init(dim,\ max\_len)}: Initializes the DTW engine.
\item
  \texttt{sentai.dtw.record\_start(label)}: Starts recording a named
  template.
\item
  \texttt{sentai.dtw.record\_add(frame)}: Adds one feature frame.
\item
  \texttt{sentai.dtw.record\_add\_imu()}: Adds the current IMU sample as
  a frame.
\item
  \texttt{sentai.dtw.record\_add\_mic()}: Adds the current microphone
  level as a frame.
\item
  \texttt{sentai.dtw.record\_end()}: Finalizes and stores the template.
\item
  \texttt{sentai.dtw.match(sequence,\ threshold)}: Matches a supplied
  sequence against stored templates.
\item
  \texttt{sentai.dtw.match\_imu(n\_frames,\ threshold,\ delay\_ms)}:
  Captures IMU frames and matches them.
\item
  \texttt{sentai.dtw.match\_mic(n\_frames,\ threshold,\ delay\_ms)}:
  Captures microphone frames and matches them.
\item
  \texttt{sentai.dtw.templates()}: Lists stored templates.
\item
  \texttt{sentai.dtw.remove(label)}: Removes a template.
\item
  \texttt{sentai.dtw.save(path)}: Saves all templates.
\item
  \texttt{sentai.dtw.load(path)}: Loads saved templates.
\item
  \texttt{sentai.dtw.info()}: Returns engine metadata.
\item
  \texttt{sentai.dtw.clear()}: Clears all templates.
\end{itemize}

\subsection{Example}\label{example-19}

\begin{verbatim}
import sentai

sentai.imu.init()
sentai.dtw.init(3, 100)
sentai.dtw.record_start('wave')
for _ in range(50):
    sentai.dtw.record_add_imu()
    sentai.rtos.sleep_ms(20)
sentai.dtw.record_end()
print(sentai.dtw.match_imu(50, 500.0))
\end{verbatim}

\section{Appendix: sentai.hmm}\label{appendix-sentai.hmm}

The \texttt{sentai.hmm} namespace implements a discrete hidden Markov
model with Baum-Welch training and Viterbi decoding. It is suitable for
compact sequence-classification tasks in which observations can be
discretized, such as simple activity recognition or discretized feature
streams extracted from sensors.

\subsection{Functions}\label{functions-20}

\begin{itemize}
\tightlist
\item
  \texttt{sentai.hmm.init(n\_states,\ n\_obs)}: Initializes the HMM.
\item
  \texttt{sentai.hmm.set\_transition(i,\ j,\ prob)}: Sets a transition
  probability.
\item
  \texttt{sentai.hmm.set\_emission(state,\ obs,\ prob)}: Sets an
  emission probability.
\item
  \texttt{sentai.hmm.set\_prior(state,\ prob)}: Sets an initial-state
  probability.
\item
  \texttt{sentai.hmm.add\_seq(obs\_list)}: Adds a training sequence.
\item
  \texttt{sentai.hmm.train(max\_iter)}: Runs Baum-Welch training.
\item
  \texttt{sentai.hmm.viterbi(obs\_list)}: Decodes the most likely state
  sequence.
\item
  \texttt{sentai.hmm.predict(obs)}: Predicts the most likely next state.
\item
  \texttt{sentai.hmm.log\_likelihood(obs\_list)}: Scores a sequence.
\item
  \texttt{sentai.hmm.from\_tpu(idx,\ n\_bins)}: Discretizes a TPU output
  tensor.
\item
  \texttt{sentai.hmm.from\_imu(n\_bins)}: Discretizes IMU magnitude.
\item
  \texttt{sentai.hmm.save(path)}: Saves the HMM model.
\item
  \texttt{sentai.hmm.load(path)}: Loads a saved model.
\item
  \texttt{sentai.hmm.info()}: Returns model metadata.
\item
  \texttt{sentai.hmm.clear()}: Clears parameters and sequences.
\end{itemize}

\subsection{Example}\label{example-20}

\begin{verbatim}
import sentai

sentai.imu.init()
sentai.hmm.init(3, 8)
seq = [sentai.hmm.from_imu(8) for _ in range(100)]
sentai.hmm.add_seq(seq)
sentai.hmm.train(20)
print(sentai.hmm.predict(sentai.hmm.from_imu(8)))
\end{verbatim}

\section{Appendix: sentai.rl}\label{appendix-sentai.rl}

The \texttt{sentai.rl} namespace collects three reinforcement-learning
mechanisms with different complexity levels: tabular Q-learning,
multi-armed bandits, and a compact DQN. It is intended for online
adaptation of thresholds, policies, or configuration choices in small
embedded experiments.

\subsection{Functions}\label{functions-21}

\begin{itemize}
\tightlist
\item
  \texttt{sentai.rl.q\_init(n\_states,\ n\_actions)}: Initializes a
  tabular Q-table.
\item
  \texttt{sentai.rl.q\_update(s,\ a,\ r,\ s\_next,\ alpha,\ gamma)}:
  Updates one Q-value.
\item
  \texttt{sentai.rl.q\_action(s,\ epsilon)}: Selects an action with an
  epsilon-greedy policy.
\item
  \texttt{sentai.rl.q\_value(s,\ a)}: Returns one Q-value.
\item
  \texttt{sentai.rl.q\_save(path)}: Saves the Q-table.
\item
  \texttt{sentai.rl.q\_load(path)}: Loads the Q-table.
\item
  \texttt{sentai.rl.q\_clear()}: Clears the Q-table.
\item
  \texttt{sentai.rl.mab\_init(n\_arms)}: Initializes a multi-armed
  bandit.
\item
  \texttt{sentai.rl.mab\_pull(arm,\ reward)}: Reports the reward of a
  selected arm.
\item
  \texttt{sentai.rl.mab\_select(strategy)}: Selects the next arm.
\item
  \texttt{sentai.rl.mab\_stats()}: Returns per-arm statistics.
\item
  \texttt{sentai.rl.mab\_clear()}: Clears bandit statistics.
\item
  \texttt{sentai.rl.dqn\_init(state\_dim,\ n\_actions,\ hidden)}:
  Initializes the DQN.
\item
  \texttt{sentai.rl.dqn\_observe(state,\ action,\ reward,\ next\_state,\ done)}:
  Adds one replay transition.
\item
  \texttt{sentai.rl.dqn\_action(state,\ epsilon)}: Selects an action
  from the DQN.
\item
  \texttt{sentai.rl.dqn\_train(batch\_size)}: Trains on a replay batch.
\item
  \texttt{sentai.rl.dqn\_save(path)}: Saves DQN weights.
\item
  \texttt{sentai.rl.dqn\_load(path)}: Loads DQN weights.
\item
  \texttt{sentai.rl.dqn\_clear()}: Clears DQN state.
\item
  \texttt{sentai.rl.info()}: Returns a summary of all three subsystems.
\end{itemize}

\subsection{Example}\label{example-21}

\begin{verbatim}
import sentai

sentai.rl.mab_init(4)
for _ in range(20):
    arm = sentai.rl.mab_select(1)
    reward = 1.0 if arm == 0 else 0.2
    sentai.rl.mab_pull(arm, reward)
print(sentai.rl.mab_stats())
\end{verbatim}

\section{Appendix: sentai.slam}\label{appendix-sentai.slam}

The \texttt{sentai.slam} namespace implements a compact detection-based
EKF-SLAM system. It estimates a planar robot pose and maintains a
landmark map using object detections or manually supplied observations.
The interface is decoupled from the TPU, so it can consume detections
from any compatible source.

\subsection{Functions}\label{functions-22}

\begin{itemize}
\tightlist
\item
  \texttt{sentai.slam.init(fov\_h\_deg,\ img\_w,\ img\_h,\ max\_lm,\ baseline\_m)}:
  Initializes the SLAM model.
\item
  \texttt{sentai.slam.update(detections)}: Updates pose and landmarks
  from monocular detections.
\item
  \texttt{sentai.slam.update\_stereo(dets\_left,\ dets\_right)}: Updates
  from stereo detections.
\item
  \texttt{sentai.slam.observe(class\_id,\ bearing\_rad,\ range\_m)}:
  Adds a manual landmark observation.
\item
  \texttt{sentai.slam.predict(dx,\ dy,\ dtheta)}: Applies an explicit
  motion prediction.
\item
  \texttt{sentai.slam.pose()}: Returns the current estimated pose.
\item
  \texttt{sentai.slam.landmarks()}: Returns the active landmark set.
\item
  \texttt{sentai.slam.noise(sigma\_v,\ sigma\_w,\ sigma\_b,\ sigma\_r)}:
  Gets or sets noise parameters.
\item
  \texttt{sentai.slam.imu\_correct(pitch\_deg,\ roll\_deg)}: Adjusts
  bearing noise using IMU tilt.
\item
  \texttt{sentai.slam.clear()}: Resets pose and landmarks.
\item
  \texttt{sentai.slam.save(path)}: Saves the map and pose.
\item
  \texttt{sentai.slam.load(path)}: Loads the map and pose.
\item
  \texttt{sentai.slam.info()}: Returns SLAM metadata and statistics.
\end{itemize}

\subsection{Example}\label{example-22}

\begin{verbatim}
import sentai

sentai.slam.init(66.0, 320, 320)
sentai.tpu.load('/models/yolov8n_320.tflite')
sentai.camera.init()

sentai.camera.to_tensor()
sentai.tpu.invoke()
dets = sentai.tpu.detect(0.3, 0.45, 50)
sentai.slam.update(dets)
print(sentai.slam.pose(), sentai.slam.landmarks())
\end{verbatim}
